\chapter{State Estimation for the Discrete-Time Linear System}

\section{Linear State Space Models}
\label{sec:6_1}

% TODO markers: operating point, T_s, sigma values, Problem 4 noise terms.

A linear state space model in discrete time has the form

\begin{align}
    x_{k+1} &= A x_k + B u_k + G w_k, \label{eq:ss-state}\\
    y_k     &= C x_k + v_k, \label{eq:ss-meas}\\
    z_k     &= C_z x_k, \label{eq:ss-out}
\end{align}

with

\begin{equation}
  \begin{bmatrix} w_k \\ v_k \end{bmatrix} \sim N_{iid}\!\left(
  \begin{bmatrix} 0 \\ 0 \end{bmatrix},
  \begin{bmatrix} R_{ww} & R_{wv} \\ R_{vw} & R_{vv} \end{bmatrix}\right),
  \qquad R_{wv}=0 .
\end{equation}

Here $u_k = [F_1 \; F_2]^T$ are the manipulated variables, $y_k = [h_1 \; h_2]_k^T$ are the measured variables and $z_k = [h_1 \; h_2]_k^T$ are the controlled variables, while the unmeasured flows $F_3$ and $F_4$ (the disturbances) are represented through the process noise $w_k$.\\

In this section, the identified transfer function model from Problem 4 and the continuous-time linearized models from Problem 5 are converted to discrete-time linear models. All the linear models are expressed in deviation variables around the operating point presented in Section \ref{sec:5_1}. The operating point is \textbf{[TODO]} and the sampling time is $T_s=\textbf{[TODO]}$~s. Zero-order hold is assumed for the pump flows, meaning that they are piecewise constant over each sampling interval, $u(t)=u_k$ for $t_k\le t<t_{k+1}=t_k+T_s$. The same holds for the disturbances $d(t)=d_k$ in the deterministic and the ZOH models, whereas the disturbances of the SDE model act continuously.

\subsection{Problem 4 models}

The Problem 4 model is an identified transfer function matrix from pump flows to levels, given in Section \ref{sec:4_4}. First, each element in the transfer function is written in the standard form

\begin{align}
    G_{ij}(s) = \frac{b^{ij}_0 s^{n_{ij}} + b^{ij}_1 s^{n_{ij}-1} + \dots + b^{ij}_{n_{ij}}}
                     {s^{n_{ij}} + a^{ij}_1 s^{n_{ij}-1} + \dots + a^{ij}_{n_{ij}}},
    \qquad b^{ij}_0=0,
\end{align}

with $n_{11}=n_{22}=1$ and $n_{12}=n_{21}=2$. Next, each element is realized as a state space model in observer canonical form

\begin{align}
  \dot x_{ij}(t) &= A_{c,ij}\,x_{ij}(t)+B_{c,ij}\,u_j(t), \\
  y_{ij}(t) &= C_{ij}\,x_{ij}(t),
\end{align}

where

\begin{equation}
  A_{c,ij}=\begin{bmatrix}-a^{ij}_1&1&&\\-a^{ij}_2&0&\ddots&\\\vdots&&&1\\ -a^{ij}_{n_{ij}}&0&\cdots&0\end{bmatrix},\quad
  B_{c,ij}=\begin{bmatrix}b^{ij}_1\\\vdots\\b^{ij}_{n_{ij}}\end{bmatrix},\quad
  C_{ij}=\begin{bmatrix}1&0&\cdots&0\end{bmatrix}.
\end{equation}

As every element is strictly proper, with $b^{ij}_0=0$, there is no direct feedthrough ($D=0$) and $B_{c,ij}$ contains only the coefficients $b^{ij}_k$. For the four elements, this gives

\begin{align}
  G_{11}:&\quad A_{c,11} = -\frac{1}{\tau_{11}},\quad
                B_{c,11} = \frac{K_{11}}{\tau_{11}},\quad
                C_{11} = 1, \\
  G_{22}:&\quad A_{c,22} = -\frac{1}{\tau_{22}},\quad
                B_{c,22} = \frac{K_{22}}{\tau_{22}},\quad
                C_{22} = 1, \\
  G_{12}:&\quad A_{c,12} = \begin{bmatrix}-a^{12}_1&1\\-a^{12}_2&0\end{bmatrix},\quad
                B_{c,12} = \begin{bmatrix}0\\ \dfrac{K_{12}}{\tau_{12,1}\tau_{12,2}}\end{bmatrix},\quad
                C_{12} = \begin{bmatrix}1&0\end{bmatrix}, \\
  G_{21}:&\quad A_{c,21} = \begin{bmatrix}-a^{21}_1&1\\-a^{21}_2&0\end{bmatrix},\quad
                B_{c,21} = \begin{bmatrix}0\\ \dfrac{K_{21}}{\tau_{21,1}\tau_{21,2}}\end{bmatrix},\quad
                C_{21} = \begin{bmatrix}1&0\end{bmatrix},
\end{align}

with

\begin{equation}
  a^{ij}_1=\frac{\tau_{ij,1}+\tau_{ij,2}}{\tau_{ij,1}\tau_{ij,2}},\qquad
  a^{ij}_2=\frac{1}{\tau_{ij,1}\tau_{ij,2}},\qquad (i,j)\in\{(1,2),(2,1)\},
\end{equation}

where $K_{ij}$, $\tau_{ij}$ and $\tau_{ij,1},\tau_{ij,2}$ are the identified values of the respective element. Subsequently, each element is integrated over one sample, since the pump flows are piecewise constant (zero-order hold). Dropping the element indices,

\begin{align}
  x_{k+1} &= Ax_k + Bu_k, \qquad
  A = e^{A_cT_s}, \quad B = \int_0^{T_s} e^{A_c\eta}\,\mathrm{d}\eta\,B_c
\end{align}

\begin{align}
    \begin{bmatrix} A & B \\ 0 & I \end{bmatrix}
  = \exp\!\left(\begin{bmatrix} A_c & B_c \\ 0 & 0 \end{bmatrix} T_s\right).
\end{align}

For the first-order elements $G_{11}$ and $G_{22}$ this yields

\begin{align}
  A_{11} &= e^{-T_s/\tau_{11}}, & B_{11} &= K_{11}\left(1-e^{-T_s/\tau_{11}}\right), \\
  A_{22} &= e^{-T_s/\tau_{22}}, & B_{22} &= K_{22}\left(1-e^{-T_s/\tau_{22}}\right),
\end{align}

and $C_{11}=C_{22}=1$. For the second-order elements $G_{12}$ and $G_{21}$, the matrices $A_{ij}$ and $B_{ij}$ are computed numerically from the same matrix exponential. Their poles follow from the continuous-time poles: $A_{c,ij}$ has the eigenvalues $-1/\tau_{ij,1}$ and $-1/\tau_{ij,2}$, and the matrix exponential maps each eigenvalue $\lambda$ to $e^{\lambda T_s}$. Hence $A_{ij}$ has the eigenvalues $e^{-T_s/\tau_{ij,1}}$ and $e^{-T_s/\tau_{ij,2}}$. \\

The output is $z_i=\sum_j G_{ij}u_j$, so the four realizations are combined into one continuous-time model

\begin{equation}
  \dot x(t)=A_cx(t)+B_cu(t),\qquad z(t)=Cx(t),\qquad
  x=\begin{bmatrix}x_{11}\\x_{12}\\x_{21}\\x_{22}\end{bmatrix},
\end{equation}

with

\begin{equation}
  A_c = \begin{bmatrix}
    A_{c,11} & 0        & 0        & 0 \\
    0        & A_{c,12} & 0        & 0 \\
    0        & 0        & A_{c,21} & 0 \\
    0        & 0        & 0        & A_{c,22}
  \end{bmatrix},
\end{equation}

\begin{equation}
  B_c = \begin{bmatrix}B_{c,11}&0\\0&B_{c,12}\\B_{c,21}&0\\0&B_{c,22}\end{bmatrix},
  \qquad
  C = \begin{bmatrix}C_{11}&C_{12}&0&0\\0&0&C_{21}&C_{22}\end{bmatrix}.
\end{equation}

Applying the zero-order-hold formula to $(A_c,B_c)$ gives

\begin{equation}
  x_{k+1}=Ax_k+Bu_k,\qquad z_k=C_zx_k,\qquad x_k\in\mathbb{R}^6,
\end{equation}

where $C_z=C$, because the measured and the controlled variables are both $h_1$ and $h_2$.

% TODO: noise and disturbance terms for the Problem 4 model (G, R_ww, R_vv),

\subsection{Problem 5 models}
\label{sec:ss61-p5}

In Problem 5, the three models from Problem 2 (deterministic, ZOH and SDE) are linearized around the operating point into continuous-time models (see Section \ref{sec:5_1}),

\begin{equation}
  \dot x(t) = A_cx(t) + B_cu(t) + E_cd(t), \qquad
  y(t) = Cx(t), \qquad z(t) = C_zx(t),
\end{equation}

where $x=[m_1\;m_2\;m_3\;m_4]^T$ are the masses in the four tanks, $d=[F_3\;F_4]^T$ are the unmeasured flows, and all variables are deviation variables. As for the Problem 4 model, the pump flows are controlled once per sample, so zero-order hold is used for the discretization. Since only the disturbance term differs between the models, the same solution over one interval can be used for all three models and has the form

\begin{align}
  x_{k+1} = e^{A_cT_s}x_k + \left(\int_0^{T_s} e^{A_c\eta}\,\mathrm{d}\eta\right)\left(B_cu_k + E_cd_k\right) = Ax_k + Bu_k + Ed_k,
\end{align}

where

\begin{align}
  A = e^{A_cT_s}, \quad
  B = \int_0^{T_s} e^{A_c\eta}\,\mathrm{d}\eta\;B_c, \quad
  E = \int_0^{T_s} e^{A_c\eta}\,\mathrm{d}\eta\;E_c,
\end{align}

and all three come from one matrix exponential,

\begin{align}
  \exp\!\left(\begin{bmatrix}
    A_c & B_c & E_c \\
    0 & 0 & 0 \\
    0 & 0 & 0
  \end{bmatrix} T_s\right) = \begin{bmatrix}
    A & B & E \\
    0 & I & 0 \\
    0 & 0 & I
  \end{bmatrix}.
\end{align}

The three disturbance terms are then discretized for each model. First, for the deterministic model the disturbances stay at their mean, $d_k=d_s$, which gives $d_k=0$ in deviation variables. The discrete linear model for the deterministic case is then

\begin{align}
  x_{k+1} = Ax_k + Bu_k, \qquad G = 0 .
\end{align}

Subsequently, for the ZOH model the flows $F_3$ and $F_4$ are over each sample defined as a random value around their mean, $d_k = w_k \sim N(0,Q_d)$ with $Q_d = \mathrm{diag}(\sigma^2_{F_3}, \sigma^2_{F_4})$. The full model is then

\begin{align}
  x_{k+1} = Ax_k + Bu_k + Ew_k, \qquad G = E, \qquad R_{ww} = Q_d,
\end{align}

with $\sigma_{F_3}=\sigma_{F_4}=\textbf{[TODO]}$~cm$^3$/s. Lastly, for the SDE model the disturbances are modeled using a Wiener process, adding directly to the mass balances of tank 3 and tank 4,

\begin{align}
  \mathrm{d}x = (A_cx + B_cu)\,\mathrm{d}t + \Sigma\,\mathrm{d}W, \qquad
  \Sigma = \rho D\begin{bmatrix} s_3 & 0 \\ 0 & s_4 \end{bmatrix}
         = E_c \begin{bmatrix} s_3 & 0 \\ 0 & s_4 \end{bmatrix},
\end{align}

with $s_3=s_4=\textbf{[TODO]}$~cm$^3$/s$^{1/2}$ and $D$ selecting tanks 3 and 4. Here, the mean disturbance is zero in the deviation variables, and $\Sigma$ does not change during the linearization, as it does not depend on the state. The full solution is therefore the same as before, but now with an added stochastic integral,

\begin{align}
  x_{k+1} = Ax_k + Bu_k + w_k, \qquad
  w_k = \int_{t_k}^{t_{k+1}} e^{A_c(t_{k+1}-\tau)}\,\Sigma\,\mathrm{d}W(\tau),
\end{align}

where $w_k$ is Gaussian with zero mean. By the It\^o isometry, and by substituting $\eta = t_{k+1}-\tau$, the covariance is

\begin{align}
  \overline{Q} = \int_0^{T_s} e^{A_c\eta}\,\Sigma\Sigma'\,e^{A_c'\eta}\,\mathrm{d}\eta,
  \qquad \Sigma\Sigma' = \rho^2 \begin{bmatrix}
    0 & 0 & 0 & 0 \\
    0 & 0 & 0 & 0 \\
    0 & 0 & s_3^2 & 0 \\
    0 & 0 & 0 & s_4^2
  \end{bmatrix}.
\end{align}

The integral is computed using one matrix exponential,

\begin{align}
  \exp\!\left(\begin{bmatrix}
    -A_c & \Sigma\Sigma' \\
    0 & A_c'
  \end{bmatrix} T_s\right) = \begin{bmatrix}
    \Phi_{11} & \Phi_{12} \\
    0 & \Phi_{22}
  \end{bmatrix}, \qquad A = \Phi_{22}', \quad \overline{Q} = \Phi_{22}'\Phi_{12},
\end{align}

meaning that $G=I$ and $R_{ww} = \overline{Q}$. Finally, the measurements are the levels in tank 1 and tank 2, which are measured with independent Gaussian noise and can therefore be expressed as

\begin{align}
  y_k = Cx_k + v_k, \qquad v_k \sim N(0,R_{vv}), \qquad R_{vv} = \sigma^2_v I, \qquad R_{wv} = 0,
\end{align}

where $C=C_z$ and $\sigma_v=\textbf{[TODO]}$~cm. The sensor noise is independent of the disturbance flows, so the cross-covariance is zero, and this holds for all three models from Problem 5.

% Chapter 4 dealt with nonlinear models, and chapter 5 attempted to linearize them.

% The chapter 4 model can therefore be modelled by an EKF, and the chapter 5 model as a standard linear Kalman Filter.

% In all cases, the measurements remain h1 and h2, wheras the unknowns are [m1,m2,m3,m4,f1,f2,f3,f4]
% We can treat the sensor inputs to f1 and f2 as noisy measurements, making it a virtual sensor.
% This gives the observable:

% y = [h1,h2,f1',f2']
% x = [m1,m2,m3,m4,f1,f2,f3,f4]

% From here, as the system in chapter 5 is linear, the continous-discrete transition is made by matrix exponentials at t + dt.

% The system transition then becomes:

% x = A*x + Gw_k
% y = C*x + v_k

% Where G is a mixing matrix for the noise, w_k and v_k are both stochastic noise terms, A is the state transition matrix and C is the observability matrix.

% Relating this to previos sections, where x = Ax + Bu
% Since the controls are not direct inputs, it becomes

% A^* = [[A,B][0,I]]
% C* = [C,0]

% For the nonlinear case, the system can still be modeled as a system of ODEs, the EKF then represents a step of a multivariate Euler integrator based on the Jacobian of the nonlinear STM.

\section{Dynamic and Static Kalman filters}

In this section, a static and a dynamic Kalman filter are designed for the identified linear models form Problem 4 and Problem 5. The static Kalman filter uses a constant gain from the stationary solution, while the dynamic Kalman filter recomputes the gain at every sample. For this case, the unknown disturbance flows $F_3$ and $F_4$ are implemented as stochastic variables without a step change. \\

The Kalman filters is designed for the discrete-time models derived in the previous Section \ref{sec:6_1}

\begin{align}
    x_{k+1} = Ax_k + Bu_k + Gw_k, \quad y_k = C_kx_k + v_k, \quad z_k = C_zx_k
\end{align}

with the initial state $x_0 = N(\overline{x}_0, P_0)$ and the noise $w_k \sim N(0,Q)$ and $v_k \sim N(0, R)$. Since the process noise is independent of each other and the starting point $x_0$, $S = cov(w_k, v_k) = 0$. Here $y = [h_1, h_2]^T$ are the measured levels and $R = \sigma^2_v I$. Relating the notions to the previous Section \ref{sec:6_1}, the same covariances where called $R_{ww} = Q$, $R_{vv} = R$ and $R_{wv} = S$. Moving in to the individual models from Problem 4 and Problem 5, the deterministic model form Problem 5 has $G = 0$ and therefore no process noise to estimate, as a consequence there is no filter needed to be designed for this model. The rest of the models where Kalman filter is designed for is presented in Table \ref{tab:6_2}.

\begin{table}[]
    \centering
    \begin{tabular}{c|c|c}
        Model & $G$ & $Q$ \\
        \hline
        Problem 5 - ZOH & $E$ & diag($\sigma^2_{F_3}, \sigma^2_{F_4}$) \\
        Problem 5 - SDE & $I$ & $\overline{Q}$ \\
        Problem 2 - ZOH & $E_B$ & diag($\sigma^2_{F_3}, \sigma^2_{F_4}$) \\
        Problem 2 - SDE & $I$ & $\overline{Q}_B$
    \end{tabular}
    \caption{Models that Kalman filter is designed for}
    \label{tab:6_2}
\end{table}


\subsection{Disturbance models}

For the Problem 5 models the disturbance models are modeled form physics, meaning that the unknown flows $F_3$ and $F_4$ enters tank 3 and tank 4 respectively, either as a random value held constant over the sample (ZOH) or as a Wiener process directly into the mass balances (SDE). In contrast, for Problem 4 the pump floes are varied in step responses with the disturbance enrty being an assumption and no distirbance path is therefore defined. This means that for the Problem 4 models, the flow $F_3$ enters tank 3 with a gain = 1, whereas $F_2$ then enters tank 2 with a gain $= (1 - \gamma_2$ leading to a transfer function for $F_3 \rightarrow h_1$ being $G_{12}(s)/(1-\gamma_2)$. Same logic can be applied to $F_4 \rightarrow h_2$ yielding the transfer function $G_{21}(s)/(1-\gamma_1)$. This means that 

\begin{align}
    B_d = \begin{bmatrix}
        0 & 0 \\
        \frac{1}{1-\gamma_2}B_{c,12} & 0 \\
        0 & \frac{1}{1-\gamma_1}B_{c,21} \\
        0 & 0 
    \end{bmatrix}
\end{align}

where a zero-order hold of $[B_e \quad B_d]$ gives $B$ and also the discrete disturbance matrix $E_B$. However, for the SDE model $\Sigma_B = B_d \; \mathrm{diag}(s_3, s_4)$ is used leading to $G = I$ and $Q = \overline{Q}_B$.\\

In all cases, the disturbances are modeled as white noise with zero-mean as is the assumption for a case without step changes. The ZOH models and SDE models do not have the same disturbance intensity as the ZOH model with a standard deviation of  $\sigma_F$ adds a variance of $\sigma^2_F T^2_s$, while the SDE model adds a variance of $s^2T_s$.

\subsection{Dynamic Kalman Filter}

For the model defined above, using one-step prediction where $\hat{x}_{k|k-1}$ is the estimate of $x_k$ based on the measurements up to sample $k-1$ (the prediction), with error covariance $P_{k|k-1}$ available at sample $k$. In addition, $x_{k|k}$ and $P_{k|k}$ are the estimate and its covariance after the measurement $y_k$ has been used (the filtered estimate. The dynamic Kalman filter will then alter between doing a measurement update and a time update.\\

Regarding the measurement update, the difference between the measured levels and the predicted levels is defined as the innovation $e_k$ with $R_{e,k}$ as its covariance and consist of the uncertainty of the predicted levels and the sensor noise

\begin{align}
    e_k = C \hat{x}_{k|k-1}, \quad R_{e,k} = CP_{k|k-1}C' + R
\end{align}

The gains for the state and the process noise can be expressed as

\begin{align}
    K_{fx,k} = P_{k|k-1}C'R_{e,k}^{-1}, \quad K_{fw,k} = SR_{e,k}^{-1}
\end{align}

which subsequently yields the filtered estimates and their covariances

\begin{align}
    \hat{x}_{k|k} &= \hat{x}_{k|k-1} + K_{fx,k}e_k, &\quad P_{k|k} = P_{k|k-1} - K_{fx,k}R_{e,k}K'_{fx,k}, \\
    \hat{w}_{k|k} &= K_{fw,k}e_k, &\quad Q_{k|k} = Q_k - K_{fw,k}R_{e,k}K'_{fw,k}
\end{align}

In the time update, the estimate is moved one sample forward with the model using the known input $u_k$

\begin{align}
    \hat{x}_{k+1|1} &= A\hat{x}_{k|k} + Bu_k + G\hat{w}_{k|k} \\
    \hat{P}_{k+1|k} &= AP_{k|k}A' + GQ_{k|k}G' - A_kK_{fx,k}S'G' - GSK'_{fx,k}A'_k
\end{align}

The prediction pushes the filtered estimate through the model and then adds the effect of the known pump flows. The last two terms of the covariance $P_{k+1|k}$ are the correction for the process noise and sensor noise when they are correlated. However, in this case the two noise terms are independent of each other, $S = 0$. This means that $K_{fw,k} = 0$, $\hat{w}_{k|k} = 0$ and $Q_{k|k} = Q$ reducing the time update to

\begin{align}
    \hat{x}_{k+1|k} = A\hat{x}_{k|k} + Bu_k, \quad P_{k+1|k} = AP_{k|k}A' + GQG'
\end{align}

The filter is initialized with the mean and covariance of the initial state, $\hat x_{0|-1}=\bar x_0=0$ and $P_{0|-1}=P_0$, with $P_0=$ [TODO].

\subsection{Static Kalman Filter}

The static Kalman filter uses the limit of P of the prediction covariance $P_{k+1|k}$, meaning that it assumes that $P_{k+1|k} \rightarrow P$ when $k \rightarrow \infty$. It then satisfies the discrete algebraic riccati equation

\begin{align}
    P = APA' + GQG' - (APC' + GS)(CPC'+R)^{-1}(APC' + GS)'
\end{align}

which for $S = 0$ the yields

\begin{align}
    P = APA' + GQG' - APC'(CPC'+R)^{-1}(APC' + GS)'CPA'
\end{align}

The constant gain and covariance are expressed as

\begin{align}
    R_e = CPC' + R, \quad K_{fx} = PC'R_e^{-1}, \quad P_f = P- K_{fx}R_eK'_{fx}
\end{align}

and the measurement and time updates defined for the dynamic Kalman filter are then used with these constant matrices. Since the covariances don't depend on the measured data, the $K_f$ is computed only once before the simulation.